\scanpage{31}
\mspage{18}

suffit de prouver $u_{I'_0,I_0}\in C$ resp. $TC$. Mais c’est là une assertion relative à l’inclusion \deleted{de \unsure{catégories}} \inserted{d’ens. ordonnés} $I'_0\subset I_0$ \uline{finis}. On est donc réduit à prouver le lemme suivant, où $I_0$ devient $I$ :

\uline{Lemme} : Soit $I$ un ens. ordonné fini, \deleted{$\unsure{I'}$ un ouvert de $I$,} $u_\bullet:X_\bullet\mathbin{\unsure{-|\to}}Y$ une flèche dans $\unsure{\mathcal M|\mathscr M}(I^\circ)$, où $\mathcal M$ est une catégorie munie d’un ens. $C\subset\mathrm{Fl}(\mathcal M)$ de flèches, tel que $C$ soit stable par coch. de base et par \inserted{\unsure{(*)}} composition, et contienne les iso.

On suppose $u_\bullet\in C_I$ i.e. que pour $i\in I$, la flèche $u_{I_{<i},I_{/i}}$ est dans $C$. Alors pour toute partie ouverte $I'$ de $I$, la flèche $u_{I',I}$ est dans $C$.

% slanted
\margin{\uline{NB} Bien sûr, $u_{J,K}\in C$ si $J\subset K$ sont deux ouverts comme on voit en appliquant le lemme à $K$ \unsure{ald} $I$, $\unsure{\to|\checkmark}$}

\inserted{\uline{\unsure{Cor.}}} Si de plus $W$ \inserted{$\subset\mathrm{Fl}\,\mathcal M$} est un localiseur modérément \unsure{ramifié|saturé}, tel que \deleted{\unsure{alors}} $TC\overset{\text{déf}}{=}W\cap C$ soit stable par coch. de base, \deleted{\unsure{alors est}} (il l’est de ttes façons par composition, et contient les iso), alors on a

\[
 C_I\cap W_{I^\circ}=(TC)_I,
\]

donc si $u_\bullet\in C_I$ (cf ci-dessus) et les $u_i\in W$, alors $u_\bullet\in(TC)_{\unsure{I}}$, donc \deleted{\unsure{par}} appliquant ce qui est dit pour $C$ à $TC$, on a $u_{I',I}\in TC$ pour \unsure{ttes|tte}
% archive sheet number: 30

